\begin{abstract}
Type error localization in Hindley--Milner systems can be formulated as a
minimum-weight Minimum Correction Subset (MWMCS) problem over unsatisfiable
type-equality constraints. We empirically study a lazy MUS-guided greedy
hitting-set approximation against an exact Z3 Optimize baseline on 300
synthetic instances spanning three structural families: single planted
conflicts, overlapping conflicts on a shared type variable, and constructor-level
function--list mismatches. Across all instances, the greedy algorithm matches
the exact optimum (approximation ratio uniformly $1.0$), and brute-force
enumeration independently validates the Z3 baseline on small instances. We
do not interpret this as evidence that greedy is generally optimal; rather,
we identify the structural reasons each generator family fails to exercise
the standard worst-case behaviour of greedy hitting set. Runtime analysis
shows that lazy greedy overhead is driven by the number of MUSes discovered
rather than the total constraint count, producing a stable $2$--$3\times$
slowdown on the dataset with overlapping conflicts. The main contribution
is a diagnostic characterisation of when synthetic HM-style generators do
and do not produce approximation gaps, which motivates future work on
adversarial generators inspired by Set Cover hardness templates.
\end{abstract}